\chapter{戦略の分類}
\label{ch:strategies}

\section{業界の分類}

ヘッジファンドの成績のデータを集める会社は、ファンドを戦略で分類している。代表的な分け方は、次の四つの大分類である（データ会社 HFR の分類にほぼ沿う）。

\begin{table}[H]
\centering\small
\caption{戦略の大分類と主な小分類}
\label{tab:taxonomy}
\begin{tabular}{@{}P{0.17\linewidth}P{0.37\linewidth}P{0.38\linewidth}@{}}
\toprule
大分類 & 主な小分類 & 一言で \\
\midrule
エクイティ・ヘッジ & ロング・ショート株式、マーケット・ニュートラル、クオンツ株式、セクター特化、ショート・バイアス & 株の銘柄選択。買いと空売りを組み合わせる。 \\
イベント・ドリブン & 合併裁定、ディストレスト（破綻・再建）、アクティビスト、特別な状況（分割・自社株買い・指数の入れ替え） & 企業の出来事の結果に賭ける。 \\
マクロ & グローバル・マクロ（裁量）、システマティック・マクロ、マネージド・フューチャーズ（CTA）、通貨 & 金利・通貨・商品・株価指数の方向に賭ける。 \\
レラティブ・バリュー & 債券の裁定、転換社債の裁定、ボラティリティの裁定、ベーシス取引、統計的裁定 & 似た資産の価格の歪みをとる。方向のリスクは打ち消す。 \\
\bottomrule
\end{tabular}
\end{table}

ほかに、これらを組み合わせるマルチストラテジー、複数のファンドに投資するファンド・オブ・ファンズがある。

\section{各戦略の中身}

ここでは、それぞれの戦略が「何を買い、何を売り、どれだけ持ち、なぜ儲かると考えられているか」を短くまとめる。研究の結果は第 II 部で詳しく扱う。

\subsection{ロング・ショート株式}

割安・好業績と判断した株を買い、割高・業績悪化と判断した株を空売りする。買いと売りの金額の差（ネット）を、資産の 0〜60\% 程度のプラスにすることが多く、市場が上がれば利益が出やすい。判断は、企業の調査（裁量）か、ファクターと統計のモデル（クオンツ）による。

\subsection{マーケット・ニュートラル}

買いと売りの金額をそろえ（ドル・ニュートラル）、さらに市場に対する感応度 $\beta$ もゼロにする（ベータ・ニュートラル）。業種や規模・スタイルのファクターへの感応度も打ち消すことが多い。残るのは、銘柄の相対的な良し悪しの予測からの収益だけである。クオンツ株式の多くはこの形で、数百〜数千の銘柄を持ち、1 銘柄の重みは小さい。

\subsection{統計的裁定}

似た値動きをする銘柄の組（ペア）や、市場全体のファクターからの「ずれ」が、元に戻ることに賭ける。保有期間は数分〜数週間で、売買の回数が多い。ずれは、一時的な需給の偏り（大口の売買）から生まれることが多く、収益は「流動性を供給した対価」と解釈される。

\subsection{トレンドフォロー（マネージド・フューチャーズ）}

株価指数・債券・通貨・商品の先物を、最近上がっていれば買い、下がっていれば売る。数十〜百以上の先物市場に分散し、保有期間は数週間〜数か月。株式の大きな下げ相場（2008 年など）でも利益を上げたことがあり、「危機のときの保険」として持たれる。

\subsection{グローバル・マクロ}

経済の見通し（成長、インフレ、金融政策、政治）から、金利・通貨・株価指数・商品の方向に賭ける。裁量のものが多いが、指標とルールで決めるシステマティック・マクロもある。代表的なルールの一つが \textbf{キャリー}：金利の高い通貨を買い低い通貨を売る、利回り曲線の形から有利な年限を持つ、などである。

\subsection{イベント・ドリブン}

\begin{itemize}
\item \textbf{合併裁定}：買収の発表後、買収される会社の株価は、買収の価格より少し低い。買収が成立すればその差が利益になる。成立しないリスクを引き受ける対価である。
\item \textbf{ディストレスト}：破綻や再建の途中の会社の債券・株を安く買い、再建後の価値に賭ける。
\item \textbf{アクティビスト}：株を大量に買い、経営に提案（自社株買い、事業の売却、経営陣の交代）をして価値を上げる。
\item \textbf{決算・指数のイベント}：決算発表の後の株価の遅れた反応（決算発表後のドリフト）、指数の銘柄の入れ替えに伴う機械的な売買など。
\end{itemize}

\subsection{レラティブ・バリュー}

\begin{itemize}
\item \textbf{債券の裁定}：ほぼ同じ国債（発行されたばかりの銘柄と、少し前に発行された銘柄）の利回りの差、国債と金利スワップの差などをとる。差は小さいので、大きなレバレッジをかける。LTCM の主な戦略だった。
\item \textbf{転換社債の裁定}：株に転換できる社債を買い、その株を空売りして株価のリスクを打ち消す。転換社債が割安に発行されやすいことから収益を得る。
\item \textbf{ベーシス取引}：現物と先物の価格の差（ベーシス）をとる。米国債の現物を買って先物を売る取引が代表で、大きなレバレッジを使う。
\end{itemize}

\subsection{ボラティリティの取引}

オプションの価格に含まれる「予想の変動率」（インプライド・ボラティリティ）は、平均すると、後で実際に起きた変動率（実現ボラティリティ）より高い。そこでオプションを売り続けると、平均的には利益が出る（ボラティリティ・リスク・プレミアム）。ただし、急落のときに大きく損をする。保険を売る商売に似ている。

\subsection{マーケットメイク・高頻度取引}

売りと買いの両方の注文を常に出し、その価格の差（スプレッド）をとる。保有期間は秒以下〜数時間。収益の源は、スプレッドと取引所の報奨金で、損の源は、情報を持った相手との取引（逆選択）と、持った在庫の値下がりである。速さと注文の仕組みの理解が勝負で、多くは自己勘定の取引会社が行う（\chapref{ch:mm}）。

\section{収益の源泉で分ける}

戦略を、「なぜ儲かるのか」で分けると、研究の整理がしやすい。

\begin{table}[H]
\centering\small
\caption{収益の源泉の分類}
\label{tab:sources}
\begin{tabular}{@{}P{0.2\linewidth}P{0.42\linewidth}P{0.3\linewidth}@{}}
\toprule
源泉 & 中身 & 例 \\
\midrule
リスク・プレミアム & 投資家が嫌うリスク（景気の悪いときに損をする、急落する）を引き受ける対価。合理的な市場でも残る。 & 株式のリスク・プレミアム、キャリー、ボラティリティの売り、合併裁定 \\
行動のバイアス & 投資家の判断の偏り（反応の遅れ、過剰な反応、注目の偏り）から生じる価格の歪み。気づく人が増えると小さくなる。 & モメンタム、決算発表後のドリフト、割安株の一部 \\
流動性の供給 & 急いで売買したい人の反対に立つ対価。 & 統計的裁定、短期リバーサル、マーケットメイク \\
構造・制約 & 規制・約款・指数のルールなど、利益以外の理由での売買。 & 指数の入れ替え、レバレッジを使えない投資家の行動（低ベータ） \\
情報・分析 & ほかの人より早く・正しく情報を処理する。 & 企業調査、代替データ、テキストの分析 \\
速さ・技術 & 注文を速く出し、速く取り消す。 & 高頻度取引、遅延の裁定 \\
\bottomrule
\end{tabular}
\end{table}

この区別は、戦略の将来を考えるうえで重要である。リスク・プレミアムは、リスクを嫌う投資家がいる限り残ると期待できるが、急落の時期を耐えなければならない。行動のバイアスによる歪みは、研究が公表され、資金が集まると小さくなりやすい（第 II 部、第 \ref{ch:failures} 章）。

\section{保有期間と容量}

戦略は、保有期間と、運用できる資金の量（容量、キャパシティ）でも大きく違う（\figref{fig:horizon}）。

\begin{figure}[H]
\centering
\begin{tikzpicture}[x=13mm,y=10mm]
\draw[arr] (0,0)--(10.3,0) node[below left,font=\footnotesize]{保有期間};
\draw[arr] (0,0)--(0,5.4) node[above,font=\footnotesize]{容量（運用できる資金）};
\foreach \x/\l in {0.5/ミリ秒,2/秒〜分,3.6/日,5.2/週,6.8/月,8.4/年} \node[lab,below] at (\x,0) {\l};
\node[lab,left] at (0,0.6) {小};
\node[lab,left] at (0,4.8) {大};
\node[box=c2,text width=19mm,minimum height=7mm] at (0.9,0.8) {マーケット\\メイク・HFT};
\node[box=c1,text width=19mm,minimum height=7mm] at (3.1,1.7) {統計的裁定};
\node[box=c4,text width=19mm,minimum height=7mm] at (4.6,2.6) {イベント\\（決算・合併）};
\node[box=c7,text width=19mm,minimum height=7mm] at (6.3,3.3) {ファクター\\株式};
\node[box=c3,text width=19mm,minimum height=7mm] at (7.4,4.6) {トレンド\\フォロー};
\node[box=ink2,text width=19mm,minimum height=7mm] at (9.1,4.0) {マクロ・\\キャリー};
\end{tikzpicture}
\caption{戦略の保有期間と容量の大まかな関係（概念図。位置は目安）。短い戦略ほど 1 回の売買の利益が小さく、取引コストに敏感で、大きな資金を入れると利益が消える。}
\label{fig:horizon}
\end{figure}

短い戦略は、1 年に何千回も独立に近い賭けができるので、シャープレシオ（\chapref{ch:perf}）は高くなりやすい。しかし、1 回の利益が小さく、売買で価格を動かしてしまうので、大きな資金は入れられない。長い戦略は、大きな資金を入れられるが、賭けの回数が少なく、シャープレシオは低くなりやすい。高頻度の取引会社が外部の資金をあまり集めず、マクロや大型のクオンツ運用会社が巨額の資金を集めるのは、この構造による。

\section{クオンツ運用の工程}
\label{sec:pipeline}

システマティックな運用は、どの戦略でも、だいたい同じ工程で作られる（\figref{fig:pipeline}）。第 I 部の残りの章は、この工程の順に並んでいる。

\begin{figure}[H]
\centering
\begin{tikzpicture}[node distance=5mm]
\node[box=ink2,text width=20mm] (d) {データ\\{\scriptsize 価格・財務・テキスト}};
\node[box=c1,text width=20mm,right=of d] (s) {シグナル\\{\scriptsize 第 5 章}};
\node[box=c1,text width=20mm,right=of s] (p) {予測\\{\scriptsize 第 6 章}};
\node[box=c7,text width=20mm,right=of p] (pc) {ポートフォリオ\\{\scriptsize 第 7 章}};
\node[box=c2,text width=20mm,right=of pc] (ex) {執行\\{\scriptsize 第 9 章}};
\node[box=c4,text width=24mm,below=9mm of p] (r) {リスクモデル\\{\scriptsize 第 8 章}};
\node[box=c4,text width=24mm,below=9mm of pc] (c) {取引コストのモデル\\{\scriptsize 第 9 章}};
\node[box=c3,text width=60mm,below=9mm of $(r)!0.5!(c)$] (bt) {検証（バックテスト・過学習の補正）\\{\scriptsize 第 10 章}};
\draw[arr] (d)--(s); \draw[arr] (s)--(p); \draw[arr] (p)--(pc); \draw[arr] (pc)--(ex);
\draw[arr] (r)--(pc); \draw[arr] (c)--(pc);
\draw[arr,dashed] (bt.west) -| (d.south);
\draw[arr,dashed] (ex.south) |- (bt.east);
\end{tikzpicture}
\caption{クオンツ運用の工程。成績の測り方（第 3 章）とファクター（第 4 章）は、すべての工程の共通の言葉になる。}
\label{fig:pipeline}
\end{figure}

\begin{projbox}
LLM-BLM の工程を当てはめると、データ＝決算発表文と株価、シグナル＝LLM・FinBERT・辞書の点数、予測＝リッジ回帰で 2 日の反応に結びつける、ポートフォリオ＝Black--Litterman と TE の上限つきの二次計画、執行＝月末の終値で売買（片道 0.3\% の費用）である。研究で弱かったのは、シグナルの「持続」（情報が当日に株価に入りきる）と、予測の大きさに比べたコストだった。プロの工程との差は、第 II 部の最後（第 \ref{ch:implications} 章）でまとめる。
\end{projbox}
